% Author: Izaak Neutelings (August 2026)
% Inspired by:
%   https://en.wikipedia.org/wiki/Confidence_interval
%   https://en.wikipedia.org/wiki/File:Normal_distribution_50%25_CI_illustration.svg
%   https://www.youtube.com/watch?v=-WVJ1o7y954 (DataMListic)
\documentclass[border=3pt,tikz]{standalone}
\usetikzlibrary{arrows.meta} % for arrow head size

% random number generator normal ~ N(mu,sigma)
% using Box-Muller transform: u1,u2 ~ U(0,1) -> z ~ N(0,1)
\usepackage{pgf}
\pgfmathdeclarefunction{normrand}{2}{%
  \pgfmathparse{#1 + #2*sqrt(-2*ln(rnd))*cos(2*pi*rnd*180/pi)}%
}

% COLORS
\colorlet{colCI}{blue!90!black}
\colorlet{colCIMissed}{red}
\colorlet{colTrue}{orange!95!black}
\colorlet{colData}{blue!60!red!80!black}

% STYLES
\tikzset{
  >={Latex[length=2,width=2]}, % for LaTeX arrow head
  CI/.style={colCI,line width=0.70},
  CIMissed/.style={CI,colCIMissed},
  true/.style={colTrue,line width=0.85},
}

\begin{document}

% CONFIDENCE INTERVAL
\foreach \Nsamples in {25,50}{%
\begin{tikzpicture}
  
  % SETTINGS
  %\def\Nsamples{25}    % number of samples/experiments
  \def\Npoints{20}     % number of data points per sample
  \def\zstar{1.959964} % z critical value for 95%
  \def\xmu{0}          % mean of normal used to generate x
  \def\xsigma{1.8}     % std. dev. of normal used to generate x
  \def\hGauss{2.3}     % peak height
  \def\hCI{0.3}        % height of each interval
  \def\yoffset{0.09}   % y offset of data points
  \def\hbarCI{0.08}    % barsize of the interval
  \pgfmathsetmacro\xmin{\xmu-3.0*\xsigma}
  \pgfmathsetmacro\xmax{\xmu+3.1*\xsigma}
  \pgfmathsetmacro\ymin{-(\Nsamples+0.2)*\hCI}
  \pgfmathsetmacro\ymax{1.2*\hGauss}
  %\pgfmathsetseed{1959964} % seed for reproducibility
  \pgfmathsetseed{12345} % seed for reproducibility
  
  % GAUSSIAN CURVE above the sample rows
  \clip (\xmin,\ymin) rectangle (\xmax,\ymax); % prevent outliers
  \draw[-{Latex[length=3,width=3]},thick,black]
    (\xmin,0) -- (\xmax,0); % axis
  \draw[true,smooth,samples=100,domain={\xmin}:{\xmax-0.1*\xsigma}]
    plot (\x,{\hGauss*exp(-(\x-\xmu)^2/(2*\xsigma*\xsigma))});
  
  % TRUE MEAN LINE
  \draw[true,dashed,line width=0.7]
    (\xmu,\ymax) node[below=5pt,left=1pt] {true mean $\mu$}
    -- (\xmu,\ymin);
  %\draw[true,dashed,line width=0.2] % std. dev.
  %  (\xmu-\xsigma,0) --++ (0,0.7*\hGauss)
  %  (\xmu+\xsigma,0) --++ (0,0.7*\hGauss);
  
  % DATA samples & intervals
  \foreach \i [evaluate={\y=-\i*\hCI; \yData=\y+\yoffset}] in {1,...,\Nsamples}{
    \message{^^JExperiment {\i/\Nsamples} with \Npoints points}
    
    % NUMBER & AXIS
    \node[above right=0pt,scale=0.6,black!70] at (\xmin,\y) {\i};
    \draw[->,line width=0.2,black!40] (\xmin,\y) -- (\xmax,\y);
    
    % DRAW data points
    \def\sumx{0}  % sum(x) to compute mean
    \def\sumsq{0} % sum(x^2) to compute standard error
    \foreach \j [evaluate={\x=normrand(\xmu,\xsigma);}] in {1,...,\Npoints}{
      \pgfmathparse{\sumx+\x}
      \xdef\sumx{\pgfmathresult}
      \pgfmathparse{\sumsq+\x*\x}
      \xdef\sumsq{\pgfmathresult}
      \fill[colData] (\x,\yData) circle (0.09*\hCI);
    }
    
    % COMPUTE mean and 95% confidence intervaL
    \pgfmathsetmacro{\mean}{\sumx/\Npoints}
    \pgfmathsetmacro{\svar}{max((\sumsq-\Npoints*\mean*\mean)/(\Npoints-1),0)}
    \pgfmathsetmacro{\stderr}{sqrt(\svar/\Npoints)}
    \pgfmathsetmacro{\cilow}{\mean - \zstar*\stderr}
    \pgfmathsetmacro{\cihigh}{\mean + \zstar*\stderr}
    \message{^^J  sumx=\sumx, sumsq=\sumsq, mean=\mean, stderr=\stderr, 95\% CI [\cilow,\cihigh]}
    
    % CHECK if CI misses the true mean
    \pgfmathsetmacro{\missesTrue}{(\cihigh<\xmu || \cilow>\xmu) ? 1 : 0}
    \ifnum \missesTrue = 1 % CI misses true mean
      \def\CIstyle{CIMissed}
    \else % CI covers true mean
      \def\CIstyle{CI}
    \fi
    
    % DRAW 95% confidence interval
    \pgfmathsetmacro{\yCI}{\yData+\yoffset}
    \draw[\CIstyle]
      (\cilow,\yCI) -- (\cihigh,\yCI) % horizontal line
      (\cilow,\yCI-\hbarCI/2) -- (\cilow,\yCI+\hbarCI/2) % bar at end
      (\cihigh,\yCI-\hbarCI/2) -- (\cihigh,\yCI+\hbarCI/2);
    \fill[\CIstyle] (\mean,\yCI) circle (0.4*\hbarCI);
    
  }
  
\end{tikzpicture}}% close loop

\end{document}